
\def\PUDcodigo{EME119}

\def\PUDcodacad{04506.35}

\def\PUDdisciplina{Sistemas Mecânicos}

\def\PUDsemestre{S7}

\def\PUDhoras{80}

%NOVOS ---------------------------------------------------
\def\PUDhorasTeoria{60}
\def\PUDhorasPratica{20}

\def\PUDinterECA{
\hyperref[PUD:EME109]{Resistência dos Materiais I (04506.26)}

\hyperref[PUD:EME111]{Resistência dos Materiais II (04506.28)}

\hyperref[PUD:EME125]{Análise de Falhas (04506.44)}
}

\def\PUDatuaisECA{---}

\def\PUDrecursos{Projetor multimídia; Quadro branco e pincel; Aulas em laboratório.}

%----------------------------------------------------------

\def\PUDcreditos{4}

\def\PUDprerequisito{ \hyperref[PUD:EME111]{EME111} }

\def\PUDpreacad{ \hyperref[PUD:EME111]{Resistência dos Materiais II (04506.28)} }

\def\PUDobjetivos{Projetar elementos de máquinas e sistemas mecânicos;  Dimensionar estruturas mecânicas;  Calcular tensões presentes em sistemas mecânicos;  Estimar vida útil de equipamentos e estruturas mecânicas.}

\def\PUDementa{Eixos Chavetas e Acoplamentos;  Mancais de Rolamento e Deslizamento;  Engrenagens Cilíndricas, Cônicas e Helicoidais;  Molas;  Parafusos e Uniões;  Embreagens e Freios;  Estudo de Casos.}

\def\PUDconteudo{ & Eixos chavetas e acoplamentos: Cargas, potências e tensões em eixos, projeto de eixos, chavetas, estrias e volantes, ajustes e acoplamentos. 
\\ 
 & Mancais de rolamento e deslizamento: Contatos conformes e não conformes, tensões de Hertz, elementos rolantes e montagens de mancais. 
\\ 
 & Engrenagens cilíndricas, cônicas e helicoidais: Teoria do dente da engrenagem, interferência e adelgaçamento, razão de contato, trem de engrenagens e tensões em engrenagens. 
\\ 
 & Molas: Constante de mola, configurações de molas, dimensionamento de molas, molas de compressão, molas de torção e molas Belleville. 
\\ 
 & Parafusos e uniões: Parafusos de potência, tensões em roscas, pré-cargas de junções em tração, fator de rigidez da junta e fixadores em cisalhamento. 
\\ 
 & Embreagens e freios: Tipos de freios e embreagens, embreagens e freios de disco e freios de tambor. 
\\ 
 & Estudo de casos. }

\def\PUDmetodo{Aulas expositórias que podem ser teóricas e/ou práticas, onde as práticas no laboratório serão marcadas no decorrer da disciplina. Desenvolvimento e construção de projetos mecânicos em nível de protótipo e, em casos especiais, em nível industrial.}

\def\PUDavalia{A avaliação será feita com aplicação de provas teóricas e/ou práticas, além da possibilidade de inclusão de trabalhos e seminários no decorrer da disciplina. Entrega final dos projetos mecânicos desenvolvidos e construídos.}

\def\PUDbibbasica{ & \href{www.biblioteca.ifce.edu.br/index.asp?codigo_sophia=40282}{Norton}, R. L.  Projeto de máquinas: uma abordagem integrada.  4ª ed.  SPorto Alegre, RS: Bookman, 2013.  1028p.  ISBN: 9788582600221. 
\\ 
 & \href{www.biblioteca.ifce.edu.br/index.asp?codigo_sophia=23902}{Collins}, J.  Projeto Mecânico de Elementos de Máquinas: uma perspectiva de prevenção da falha.  1ª ed.  Rio de Janeiro, RJ: LTC, 2006.  740p.  ISBN: 8521614756. 
\\ 
 & \href{www.biblioteca.ifce.edu.br/index.asp?codigo_sophia=24029}{Melconian}, S.  Elementos de Máquinas.  8ª ed.  São Paulo, SP: Érica, 2007.  358p.  ISBN: 9788571947030. }

\def\PUDbibcomplementar{ & \href{www.biblioteca.ifce.edu.br/index.asp?codigo_sophia=64982}{Ashby}, M. F.  Seleção de Materiais no Projeto Mecânico.  1º ed.  Rio de Janeiro, RJ: Elsevier, 2012.  673p.  ISBN: 9788535245219.
%&  SHIGLEY, J.  E.;  MISCHKE, C.  R.;  BUDYNAS, R.  G.;  Projeto de Engenharia Mecânica.  7ª edição.  São Paulo.  Editora BOOKMAN.  2005.  960p.  ISBN 9788536305622. 
\\ 
 &  \href{www.biblioteca.ifce.edu.br/index.asp?codigo_sophia=63128}{Juvinall}, R. C.; Marshek, K. M.  Fundamentos do Projeto de Componentes de Máquinas.  5ª ed.  Rio de Janeiro, RJ: LTC, 2016.  562p.  ISBN: 9788521630098. 
\\
 & \href{www.biblioteca.ifce.edu.br/index.asp?codigo_sophia=58590}{Mott}, Robert L.  Elementos de máquina em projetos mecânicos.  E-book.  5º ed.  Pearson.  924p.  ISBN: 9788543005904.
\\
 & \href{www.biblioteca.ifce.edu.br/index.asp?codigo_sophia=61931}{Johnston Jr.}, E. R.;  Beer, F. P.;  Mecânica dos Materiais.  7ª ed.  Porto Alegre, RS: AMGH, 2015.  838p.  ISBN: 9788580554984. 
 \\
 & \href{www.biblioteca.ifce.edu.br/index.asp?codigo_sophia=59395}{Pereira}, Celso Pinto Morais.  Mecânica dos materiais avançada.  E-book.  Interciência.  434p.  ISBN: 9788571933347.}

\def\PUDautor{Venceslau Xavier de Lima Filho}

\def\PUDdata{2013-06-17}

\def\PUDrevisao{2}

\def\PUDrevisor{Venceslau Xavier de Lima Filho}

\def\PUDdatarevisao{2019-05-15}

\PUDcorpo
